\section*{NeurIPS Paper Checklist}

\small
\begin{enumerate}[leftmargin=*,itemsep=2pt,parsep=0pt,topsep=2pt]

\item {\bf Claims}
    \item[] Question: Do the main claims made in the abstract and introduction accurately reflect the paper's contributions and scope?
    \item[] Answer: \answerYes{}
    \item[] Justification: The abstract and introduction claim that \emph{in the frozen and light-adaptation regime of medical segmentation pipelines}, same-FM/same-recipe pools show higher per-case Dice-score correlation, used as a failure-dependence proxy, than the Table-1 cross-family reference pools. The primary table uses task-specific post-floor analytic pools: Kvasir 11 FMs, ACDC LV 10, BraTS foreground 10, RIGA Cup 9, and RIGA Disc 11. The scope is explicitly restricted to retrospective public-benchmark dependence audits; secondary pixel/lift readouts, released nnU-Net complements, and non-primary mechanism diagnostics are diagnostic or scope context rather than causal clinical-outcome evidence.

\item {\bf Limitations}
    \item[] Question: Does the paper discuss the limitations of the work performed by the authors?
    \item[] Answer: \answerYes{}
    \item[] Justification: The main paper gives high-level limitations in Section~\ref{sec:discussion}, and the appendix expands them with scoped protocol notes: frozen/light-adaptation only; unbundled LoRA and RIGA full-fine-tuning mechanism probes retained only as transparency notes; released CNN random-init and small ACDC full-fine-tuning provenance traces are diagnostic, not primary evidence; UNet-with-skip and task-level nnU-Net v2 complements reported but not fully case-identical dense 3D FM/full-pipeline baselines for every primary row; the unbundled 3D architecture pilot is appendix scope context only; 2D axial slices at $224{\times}224$ for the primary FM audit; MedSAM used only as a frozen vision-encoder probe; broader cross-scanner / cross-demographic shifts remain untested; clinical Dice thresholds are operational conventions and many lift thresholds are support-limited.

\item {\bf Theory assumptions and proofs}
    \item[] Answer: \answerNA{}
    \item[] Justification: The paper is empirical and benchmark-oriented; it does not include theoretical results or proofs. Section~\ref{sec:formal} introduces estimand notation but the $\bar{r}$, $\Delta$, and lift statistics are descriptive.

\item {\bf Experimental result reproducibility}
    \item[] Answer: \answerYes{}
    \item[] Justification: Section~\ref{sec:design} and Appendix~\ref{sec:adaptation} specify the public FM checkpoints, decoder architecture ($1{\times}1$-conv default plus stated sweeps), training protocol (Adam $10^{-3}$, 30 epochs for the frozen-head primary runs, batch 16, BCE for the binary primary rows), seed set ($\{42,43,44,45\}$ with $\sigma$-seeded determinism), and the actual primary split protocols: Kvasir 880/120, RIGA train = BinRushed+MESSIDOR with Magrabia (Magrabi Eye Center) source held out ($n{=}95$), ACDC public patients 001--100 with patient-level fold-0 LV-positive slices ($n{=}361$), and the BraTS 2024 2D FLAIR \texttt{seg>0} foreground derivative: per held-out subject, up to the top-10 axial slices by foreground area with at least 64 positive pixels, subject-level 80/20 split using NumPy seed 42 ($n{=}3228$). This BraTS target includes the GLI resection-cavity label and is not the official 3D BraTS WT challenge target. The supplementary artifact releases primary-task per-case Dice JSONs, merged summaries, selected aggregate diagnostic JSONs, and code; full raw data, full prediction caches, raw held-out prediction arrays, and fuller split manifests are outside the supplementary artifact.

\item {\bf Open access to data and code}
    \item[] Answer: \answerYes{}
    \item[] Justification: All source datasets are public or access-controlled research resources from their upstream providers, subject to their licences, research-use terms, or DUAs (RIGA+ via Deep Blue; ACDC via the ACDC challenge; Kvasir-SEG via Simula; BraTS 2024 adult glioma via Synapse; CVC-ClinicDB via an operational HuggingFace mirror for a no-raw-redistribution OOD diagnostic). All backbones are publicly available or documented as gated research checkpoints. The supplementary artifact releases benchmark code, estimator implementations, primary-task per-case Dice traces, merged summaries, and selected appendix result JSONs; authored code and metadata use the bundle licence, while derived scalar traces remain subject to upstream dataset/model terms. Raw medical images, full prediction caches, some held-out-task artifacts, and complete split manifests are not redistributed in the supplementary artifact.

\item {\bf Experimental setting/details}
    \item[] Answer: \answerYes{}
    \item[] Justification: Section~\ref{sec:formal} gives the formal estimands (within-/cross-family $\bar{r}$, spatial $\Delta_{\mathrm{pix}}$, threshold lift $L$) and the symmetric $4{\times}4$ seed-combo averaging rule; Section~\ref{sec:design} gives the datasets (with patient-level splits), foundation-model list, and protocol (image size, encoder-specific upstream normalisation, decoder, optimiser, loss, epochs, batch, seed control).

\item {\bf Experiment statistical significance}
    \item[] Answer: \answerYes{}
    \item[] Justification: Table~\ref{tab:main} reports paired case-level and hierarchical family/seed bootstrap 95\% CIs ($B{=}2000$). The BraTS foreground derivative is analysed at the slice level with an explicit within-subject dependence caveat, Appendix~\ref{app:subject-cluster} reports subject-level sensitivity checks, and Appendix~\ref{app:brats-ablations} explains why sparse-class NCR filtering diagnostics are not used as numeric evidence.

\item {\bf Experiments compute resources}
    \item[] Answer: \answerYes{}
    \item[] Justification: Experiments were run on NVIDIA A100-80GB, A6000, and L40S GPUs across a shared compute cluster. Feature extraction per FM per domain ranges from minutes for CNN backbones to longer ViT/BraTS jobs; cached 1$\times$1-head training per seed is lightweight. Exact wall-clock varies by completed matrix; the release reports per-job elapsed times in result JSONs where recorded rather than claiming a single exact compute total.

\item {\bf Code of ethics}
    \item[] Answer: \answerYes{}
    \item[] Justification: The research uses only de-identified public or access-controlled research medical imaging datasets (RIGA, ACDC, Kvasir-SEG, BraTS 2024) under their original research licences, terms, or DUAs. No additional patient data are collected. The paper audits retrospective failure-correlation properties of benchmark segmentation pipelines and does not propose new clinical tools.

\item {\bf Broader impacts}
    \item[] Answer: \answerYes{}
    \item[] Justification: Section~\ref{sec:discussion} discusses deployment implications: shared-pipeline model pools (same FM, same architecture family, same recipe) in frozen/light-adaptation medical AI deployments may be over-estimated as providing independent safety margins. The paper quantifies a retrospective evaluation assumption (independence among FM-pool members) and does not measure clinician decisions, diagnoses, or patient outcomes. Backbone-family diversification is observed to have lower dependence in the tested comparisons, not to be a complete clinical mitigation.

\item {\bf Safeguards}
    \item[] Answer: \answerYes{}
    \item[] Justification: The supplementary artifact does not release raw medical images, voxel labels, full prediction caches, upstream checkpoints, trained clinical services, or patient-facing models. Released artifacts are derived per-case scalar traces, aggregate summaries, metadata, hashes, and scorer code under upstream dataset/model constraints. The paper states that the audit is not a clinical safety guarantee, does not support re-identification or patient profiling, and should not be used as a standalone model-selection or deployment rule.

\item {\bf Licenses for existing assets}
    \item[] Answer: \answerYes{}
    \item[] Justification: RIGA+ (CC BY-NC 4.0 via Deep Blue; non-commercial use only; raw images and masks not redistributed), ACDC (research licence via challenge website), Kvasir-SEG (research/educational use via Simula; raw images not redistributed), BraTS 2024 (research licence/DUA via Synapse; no raw images or voxel-level labels redistributed), and CVC-ClinicDB used only through a no-raw-redistribution OOD diagnostic; we do not redistribute raw CVC images or masks, and any future raw redistribution would require checking the original CVC terms. Backbones/checkpoints: DINOv2 (Apache-2.0), MAE ViT-B/16 (CC BY-NC 4.0), DeiT-B/16 (Apache-2.0), OpenAI CLIP ViT-B/16 and ViT-L/14 weights (OpenAI CLIP license; open\_clip code MIT), BiomedCLIP checkpoint under its model-card terms rather than the open\_clip library license, ConvNeXt-T/EfficientNet-B0/ResNet-50/18 loaded via torchvision weights/code under PyTorch/torchvision terms, RETFound checkpoint (gated; CC BY-NC 4.0/research-use), MedSAM/SAM (Apache-2.0). No upstream checkpoints are redistributed. \textbf{BraTS redistribution compliance:} in the supplementary artifact we release only derived scalar traces and aggregate statistics needed to reproduce the estimators; we do \emph{not} redistribute raw BraTS images, raw voxel labels, or full prediction caches. All assets are cited.

\item {\bf New assets}
    \item[] Answer: \answerYes{}
    \item[] Justification: The paper releases (i)~a benchmark protocol (estimands + symmetric multi-seed + bootstrap estimators); (ii)~primary-task per-case Dice JSONs and merged summary JSONs for the released rows; (iii)~selected appendix/ablation result JSONs and reproduction scripts. Raw prediction masks/caches and full split manifests are not part of the supplementary artifact because upstream redistribution rights and artifact size constraints vary across datasets.

\item {\bf Crowdsourcing and research with human subjects}
    \item[] Answer: \answerNA{}
    \item[] Justification: No crowdsourcing; no new human-subjects research. We use existing public or access-controlled research datasets whose upstream governance, approvals, or access terms are handled by the dataset providers.

\item {\bf Institutional review board (IRB) approvals or equivalent for research with human subjects}
    \item[] Answer: \answerNA{}
    \item[] Justification: No new human-subjects data were collected. Upstream datasets carry their own governance, approvals, or access terms (RIGA, ACDC, Kvasir-SEG, BraTS 2024).

\item {\bf Declaration of LLM usage}
    \item[] Answer: \answerYes{}
    \item[] Justification: LLMs were used as editorial and programming assistants during manuscript and artifact preparation. They are not a component of the benchmark methodology, estimator, training pipeline, or evaluation protocol: all reported numbers come from the released scripts, logs, and JSON artifacts, and all manuscript text, code, citations, and claims were reviewed and approved by the authors.

\end{enumerate}
